% Extended Mathematical Proofs
% Supplementary Material for: Co-occurrence, Sequence and Knowledge Graph 
% with Ontology as a Second Brain for AI-LLM

\documentclass[11pt]{article}
\usepackage[utf8]{inputenc}
\usepackage{amsmath,amssymb,amsfonts,amsthm}
\usepackage[margin=1in]{geometry}
\usepackage{hyperref}

\newtheorem{theorem}{Theorem}
\newtheorem{lemma}[theorem]{Lemma}
\newtheorem{proposition}[theorem]{Proposition}
\newtheorem{corollary}[theorem]{Corollary}
\newtheorem{definition}{Definition}
\newtheorem{remark}{Remark}

\title{Extended Mathematical Proofs\\
\large Supplementary Material for Triple-Source RAG Framework}
\author{Anonymous}
\date{}

\begin{document}

\maketitle

\section{Introduction}

This document provides complete proofs for all theorems stated in the main paper. We organize proofs by topic: gating mechanism properties (Section~\ref{sec:gating}), training convergence (Section~\ref{sec:training}), complexity bounds (Section~\ref{sec:complexity}), and performance analysis (Section~\ref{sec:performance}).

\section{Gating Mechanism Properties}
\label{sec:gating}

\subsection{Preliminary Definitions}

\begin{definition}[Sigmoid Function]
The sigmoid function $\sigma: \mathbb{R} \to (0, 1)$ is defined as:
\begin{equation}
\sigma(x) = \frac{1}{1 + e^{-x}}
\end{equation}
\end{definition}

\begin{definition}[Softmax Function]
The softmax function $\text{softmax}: \mathbb{R}^n \to \Delta^{n-1}$ is defined as:
\begin{equation}
\text{softmax}(\mathbf{z})_j = \frac{e^{z_j}}{\sum_{k=1}^n e^{z_k}}
\end{equation}
where $\Delta^{n-1}$ is the $(n-1)$-dimensional probability simplex.
\end{definition}

\subsection{Theorem 1: Gating Bounds}

\begin{theorem}[Gating Bounds - Complete Statement]
\label{thm:bounds_full}
Let $g_{\text{cand}}: \mathbb{R}^d \to \mathbb{R}$ and $g_{\text{src}}: \mathbb{R}^d \to \mathbb{R}^3$ be defined as:
\begin{align}
g_{\text{cand}}(i) &= \sigma(\mathbf{w}^\top \mathbf{h}_i) \\
\mathbf{g}_{\text{src}} &= \text{softmax}\left(\frac{\mathbf{W}_g \mathbf{h}_q}{\tau}\right)
\end{align}
where $\mathbf{w} \in \mathbb{R}^d$, $\mathbf{W}_g \in \mathbb{R}^{3 \times d}$, $\tau > 0$. Then:
\begin{enumerate}
    \item $g_{\text{cand}}(i) \in (0, 1)$ for all $\mathbf{h}_i \in \mathbb{R}^d$
    \item $\mathbf{g}_{\text{src}} \in \Delta^2$ with $\sum_{s=1}^3 g_{\text{src}}(s) = 1$ and $g_{\text{src}}(s) > 0$
    \item $\text{score}(i, s) = g_{\text{cand}}(i) \times g_{\text{src}}(s) \in (0, 1)$
\end{enumerate}
\end{theorem}

\begin{proof}
We prove each part:

\textbf{Part 1:} For any $\mathbf{h}_i \in \mathbb{R}^d$, the inner product $\mathbf{w}^\top \mathbf{h}_i \in \mathbb{R}$ is a real number. We analyze the sigmoid:

\begin{itemize}
    \item $\sigma(x) > 0$ for all $x \in \mathbb{R}$ since $e^{-x} > 0$, hence $1 + e^{-x} > 1 > 0$
    \item $\sigma(x) < 1$ for all $x \in \mathbb{R}$ since $e^{-x} > 0$ implies $1 + e^{-x} > 1$
    \item Limits: $\lim_{x \to -\infty} \sigma(x) = 0$ and $\lim_{x \to +\infty} \sigma(x) = 1$
\end{itemize}

Therefore $\sigma(\mathbf{w}^\top \mathbf{h}_i) \in (0, 1)$ for any finite $\mathbf{h}_i$.

\textbf{Part 2:} For the softmax function with input $\mathbf{z} = \mathbf{W}_g \mathbf{h}_q / \tau$:

\begin{itemize}
    \item $e^{z_j} > 0$ for all $j$ since exponential is always positive
    \item $\sum_{j=1}^3 e^{z_j} > 0$ is a sum of positive terms
    \item Therefore $\text{softmax}(\mathbf{z})_j = e^{z_j} / \sum_{k} e^{z_k} > 0$
\end{itemize}

For the sum:
\begin{equation}
\sum_{j=1}^3 \text{softmax}(\mathbf{z})_j = \sum_{j=1}^3 \frac{e^{z_j}}{\sum_{k=1}^3 e^{z_k}} = \frac{\sum_{j=1}^3 e^{z_j}}{\sum_{k=1}^3 e^{z_k}} = 1
\end{equation}

\textbf{Part 3:} Since $g_{\text{cand}}(i) \in (0, 1)$ and $g_{\text{src}}(s) \in (0, 1)$:
\begin{equation}
0 < g_{\text{cand}}(i) \times g_{\text{src}}(s) < 1 \times 1 = 1
\end{equation}
Thus $\text{score}(i, s) \in (0, 1)$.
\end{proof}

\subsection{Theorem 2: Ranking Consistency}

\begin{theorem}[Ranking Consistency]
\label{thm:ranking_full}
Let $\{s_1, \ldots, s_n\}$ be scores computed by the gating mechanism for $n$ candidates. Define $\gamma = \min_{i \neq j} |s_i - s_j|$. For any perturbation $\{\delta_i\}$ with $|\delta_i| < \gamma/2$ for all $i$, the ranking induced by $\{s_i + \delta_i\}$ is identical to that induced by $\{s_i\}$.
\end{theorem}

\begin{proof}
Suppose $s_i > s_j$ for some $i \neq j$. We show $(s_i + \delta_i) > (s_j + \delta_j)$:

\begin{align}
(s_i + \delta_i) - (s_j + \delta_j) &= (s_i - s_j) + (\delta_i - \delta_j) \\
&\geq (s_i - s_j) - |\delta_i - \delta_j| \\
&\geq \gamma - (|\delta_i| + |\delta_j|) \\
&> \gamma - (\gamma/2 + \gamma/2) \\
&= 0
\end{align}

Therefore all pairwise orderings are preserved, and the ranking is identical.
\end{proof}

\subsection{Gradient Properties}

\begin{lemma}[Gating Gradients]
\label{lem:gradients}
The gating functions have well-defined gradients:
\begin{align}
\frac{\partial g_{\text{cand}}(i)}{\partial \mathbf{w}} &= \sigma'(\mathbf{w}^\top \mathbf{h}_i) \cdot \mathbf{h}_i \\
\frac{\partial g_{\text{src}}(s)}{\partial \mathbf{W}_g} &= \frac{1}{\tau} \left( g_{\text{src}}(s) \cdot \delta_{sj} - g_{\text{src}}(s) \cdot g_{\text{src}}(j) \right) \mathbf{h}_q^\top
\end{align}
where $\sigma'(x) = \sigma(x)(1 - \sigma(x))$ is the sigmoid derivative.
\end{lemma}

\begin{proof}
For the sigmoid:
\begin{equation}
\frac{d\sigma(x)}{dx} = \frac{e^{-x}}{(1+e^{-x})^2} = \sigma(x)(1-\sigma(x))
\end{equation}

By chain rule: $\frac{\partial g_{\text{cand}}(i)}{\partial \mathbf{w}} = \sigma'(\mathbf{w}^\top \mathbf{h}_i) \cdot \mathbf{h}_i$.

For softmax, the Jacobian is:
\begin{equation}
\frac{\partial \text{softmax}(\mathbf{z})_s}{\partial z_j} = \text{softmax}(\mathbf{z})_s (\delta_{sj} - \text{softmax}(\mathbf{z})_j)
\end{equation}

The chain rule with $\mathbf{z} = \mathbf{W}_g \mathbf{h}_q / \tau$ gives the result.
\end{proof}

\section{Training Convergence}
\label{sec:training}

\subsection{InfoNCE Properties}

\begin{definition}[InfoNCE Loss]
For query $q$ with positive sample $p^+$ and $N-1$ negatives $\{p^-_j\}$:
\begin{equation}
\mathcal{L}_{\text{InfoNCE}} = -\log \frac{\exp(f(q, p^+)/\tau)}{\sum_{j=1}^N \exp(f(q, p_j)/\tau)}
\end{equation}
where $f(q, p) = \cos(\mathbf{h}_q, \mathbf{h}_p)$ is cosine similarity.
\end{definition}

\begin{theorem}[InfoNCE and Mutual Information]
\label{thm:infonce_full}
Let $I(Q; P^+)$ denote the mutual information between query and positive passage distributions. Then:
\begin{equation}
I(Q; P^+) \geq \log(N) - \mathbb{E}_{q, p^+, \{p^-\}}[\mathcal{L}_{\text{InfoNCE}}]
\end{equation}
\end{theorem}

\begin{proof}
This is the standard result from van den Oord et al. (2018). The key insight is that InfoNCE is a lower bound on mutual information via the data processing inequality.

Let $p(q, p)$ be the joint distribution and $p(q)p(p)$ the product of marginals. The mutual information is:
\begin{equation}
I(Q; P) = \mathbb{E}_{p(q,p)}\left[\log \frac{p(q, p)}{p(q)p(p)}\right]
\end{equation}

The InfoNCE objective estimates the density ratio $\frac{p(p|q)}{p(p)}$ using a classifier with $N$ classes. The bound follows from the optimal Bayes classifier having accuracy $\frac{e^{I(Q;P)}}{e^{I(Q;P)} + N - 1}$.
\end{proof}

\subsection{Convergence Guarantee}

\begin{theorem}[SGD Convergence for InfoNCE]
\label{thm:convergence_full}
Let $\mathcal{L}(\theta)$ be the InfoNCE loss parameterized by $\theta$. Under the following conditions:
\begin{enumerate}
    \item $\mathcal{L}$ is $L$-Lipschitz smooth
    \item Gradients are bounded: $\|\nabla \mathcal{L}(\theta)\| \leq G$
    \item Learning rate schedule: $\eta_t = \eta_0 / \sqrt{t}$
\end{enumerate}
SGD converges to a stationary point: $\mathbb{E}[\|\nabla \mathcal{L}(\theta_T)\|^2] \leq O(1/\sqrt{T})$.
\end{theorem}

\begin{proof}
InfoNCE satisfies smoothness since it is a composition of:
\begin{itemize}
    \item Cosine similarity (bounded in $[-1, 1]$)
    \item Division by temperature $\tau > 0$
    \item Exponential (smooth)
    \item Log-sum-exp (smooth)
    \item Logarithm (smooth when argument bounded away from 0)
\end{itemize}

The denominator $\sum_j \exp(f(q, p_j)/\tau) \geq \exp(-1/\tau) > 0$ is bounded away from zero since cosine similarity is in $[-1, 1]$.

Standard SGD convergence (Bottou et al., 2018) then applies:
\begin{equation}
\min_{t \leq T} \mathbb{E}[\|\nabla \mathcal{L}(\theta_t)\|^2] \leq \frac{\mathcal{L}(\theta_0) - \mathcal{L}^*}{\sum_{t=1}^T \eta_t} + \frac{L\sigma^2 \sum_{t=1}^T \eta_t^2}{2\sum_{t=1}^T \eta_t}
\end{equation}

With $\eta_t = O(1/\sqrt{t})$, both terms are $O(1/\sqrt{T})$.
\end{proof}

\section{Complexity Analysis}
\label{sec:complexity}

\subsection{Time Complexity}

\begin{theorem}[Pipeline Time Complexity - Detailed]
\label{thm:time_full}
The complete retrieval pipeline has time complexity:
\begin{equation}
T(n, k, d, |V|) = O\left(\sqrt{n} \cdot d + k^2 \cdot d + |V| \cdot d_c + |E_q| \cdot \bar{d}\right)
\end{equation}
where:
\begin{itemize}
    \item $n$ = corpus size
    \item $k$ = candidates per source
    \item $d$ = embedding dimension
    \item $|V|$ = vocabulary size
    \item $d_c$ = co-occurrence embedding dimension
    \item $|E_q|$ = number of linked entities in query
    \item $\bar{d}$ = average degree in knowledge graph
\end{itemize}
\end{theorem}

\begin{proof}
We analyze each component:

\textbf{1. Query Encoding:} $O(d \cdot |q|)$ where $|q|$ is query length. Negligible for short queries.

\textbf{2. FAISS IVF Retrieval:} With IVF index, we probe $\sqrt{n}$ centroids and search within those partitions. Cost: $O(\sqrt{n} \cdot d)$.

\textbf{3. Co-occurrence Retrieval:} 
\begin{itemize}
    \item Naive: $O(|V| \cdot d_c)$ to scan all embeddings
    \item With index: $O(\sqrt{|V|} \cdot d_c)$
\end{itemize}

\textbf{4. KG Retrieval:}
\begin{itemize}
    \item Entity linking: $O(|q| \cdot |\mathcal{E}|)$ or $O(|q|)$ with BLINK
    \item Neighborhood expansion: $O(|E_q| \cdot \bar{d}^h)$ for $h$-hop
    \item Passage ranking: $O(|E_q| \cdot \bar{d} \cdot k)$
\end{itemize}

\textbf{5. Cross-Attention:} For $3k$ candidates with $d$-dimensional embeddings:
\begin{itemize}
    \item Attention matrix: $O((3k)^2 \cdot d) = O(k^2 \cdot d)$
    \item Three source pairs: $O(k^2 \cdot d)$ each
    \item Total: $O(k^2 \cdot d)$
\end{itemize}

\textbf{6. Gating:}
\begin{itemize}
    \item Per-candidate: $O(k \cdot d)$ for $3k$ dot products
    \item Source-level: $O(d)$ for query projection
    \item Total: $O(k \cdot d)$
\end{itemize}

\textbf{7. Top-k Selection:} Using heap: $O(3k \cdot \log k')$ where $k'$ is final count.

Combining: Cross-attention $O(k^2 \cdot d)$ and FAISS $O(\sqrt{n} \cdot d)$ dominate.
\end{proof}

\subsection{Space Complexity}

\begin{theorem}[Space Complexity - Detailed]
\label{thm:space_full}
Storage and working memory requirements:
\begin{align}
S_{\text{storage}} &= O(n \cdot d + |V|^2 + |\mathcal{T}| + d^2) \\
S_{\text{working}} &= O(k \cdot d)
\end{align}
\end{theorem}

\begin{proof}
\textbf{Storage:}
\begin{itemize}
    \item FAISS index: $O(n \cdot d)$ for $n$ embeddings
    \item Co-occurrence: $O(|V|^2)$ dense or $O(\text{nnz})$ sparse
    \item Knowledge graph: $O(|\mathcal{T}|)$ for triples
    \item Model parameters: $O(d^2)$ for attention weights
\end{itemize}

\textbf{Working Memory:}
\begin{itemize}
    \item Candidate embeddings: $O(3k \cdot d)$
    \item Attention matrices: $O(k^2)$ per layer
    \item Intermediate activations: $O(k \cdot d)$
\end{itemize}
\end{proof}

\section{Performance Analysis}
\label{sec:performance}

\subsection{Coverage Improvement}

\begin{proposition}[Coverage Improvement - Detailed]
\label{prop:coverage_full}
Let $R_A$, $R_B$, $R_C$ be recall of three sources. Under conditional independence given relevance:
\begin{equation}
R_{A \cup B \cup C} = 1 - (1 - R_A)(1 - R_B)(1 - R_C)
\end{equation}
\end{proposition}

\begin{proof}
By definition, recall is the probability that a relevant item is retrieved. Let $R$ denote the event that a relevant item is retrieved. Then:
\begin{equation}
P(R_A \cup R_B \cup R_C) = 1 - P(\overline{R_A} \cap \overline{R_B} \cap \overline{R_C})
\end{equation}

Under conditional independence:
\begin{equation}
P(\overline{R_A} \cap \overline{R_B} \cap \overline{R_C} | \text{relevant}) = (1-R_A)(1-R_B)(1-R_C)
\end{equation}

Therefore:
\begin{equation}
R_{A \cup B \cup C} = 1 - (1-R_A)(1-R_B)(1-R_C)
\end{equation}

\textbf{Numerical example:} If $R_A = R_B = R_C = 0.7$:
\begin{equation}
R_{A \cup B \cup C} = 1 - (0.3)^3 = 1 - 0.027 = 0.973
\end{equation}

Combined recall (97.3\%) substantially exceeds individual recall (70\%).
\end{proof}

\subsection{Adaptive Optimality}

\begin{proposition}[Adaptive Optimality - Detailed]
\label{prop:adaptive_full}
For query distribution $P(Q)$ with query-dependent optimal weights $\mathbf{w}^*(q)$, learned gating $\hat{\mathbf{w}}(q)$ achieves expected loss at most equal to any fixed weighting $\mathbf{w}_{\text{fixed}}$:
\begin{equation}
\mathbb{E}_q[\mathcal{L}(q, \hat{\mathbf{w}}(q))] \leq \mathbb{E}_q[\mathcal{L}(q, \mathbf{w}_{\text{fixed}})]
\end{equation}
with strict inequality when optimal weights vary across queries.
\end{proposition}

\begin{proof}
For any fixed $\mathbf{w}_{\text{fixed}}$:
\begin{align}
\mathbb{E}_q[\mathcal{L}(q, \hat{\mathbf{w}}(q))] &\leq \mathbb{E}_q[\min_\mathbf{w} \mathcal{L}(q, \mathbf{w})] \quad \text{(definition of } \hat{\mathbf{w}}\text{)} \\
&\leq \mathbb{E}_q[\mathcal{L}(q, \mathbf{w}_{\text{fixed}})] \quad \text{(min $\leq$ any choice)}
\end{align}

For strict inequality, consider two query types $q_1, q_2$ with $P(q_1) = P(q_2) = 0.5$ and:
\begin{itemize}
    \item $\mathbf{w}^*(q_1) = (1, 0, 0)$ (source A optimal)
    \item $\mathbf{w}^*(q_2) = (0, 1, 0)$ (source B optimal)
\end{itemize}

Any fixed weighting must compromise, while adaptive weighting achieves optimal loss for each query type.
\end{proof}

\section{Conclusion}

This document provides complete proofs for all theoretical claims in the main paper. The key results are:

\begin{enumerate}
    \item The gating mechanism is well-defined with bounded outputs (Theorem~\ref{thm:bounds_full})
    \item InfoNCE training maximizes a mutual information lower bound (Theorem~\ref{thm:infonce_full})
    \item SGD on InfoNCE converges to a stationary point (Theorem~\ref{thm:convergence_full})
    \item Pipeline complexity is tractable: $O(\sqrt{n} \cdot d + k^2 \cdot d)$ (Theorem~\ref{thm:time_full})
    \item Multi-source fusion improves coverage under independence (Proposition~\ref{prop:coverage_full})
    \item Adaptive gating is strictly better than fixed weighting (Proposition~\ref{prop:adaptive_full})
\end{enumerate}

\end{document}
